Par une nuit claire et sans vent, le sol d'un jardin, d'émissivité \num{\v{em}}, est à \qty{\v{Tg}}{\degreeCelsius}. Le ciel clair rayonne vers le sol comme un corps noir à une température effective de \qty{\v{Tc}}{\degreeCelsius}.

\begin{enumerate}[label=\alph*)]
    \item Déterminez la puissance nette perdue par rayonnement par chaque mètre carré de sol.
    \item Une couche de nuages arrive ; la température effective du ciel monte à \qty{\v{Tn}}{\degreeCelsius}. Déterminez la nouvelle perte nette par mètre carré.
    \item Expliquez pourquoi le gel au sol est plus fréquent par nuit claire que par nuit nuageuse, même à température de l'air égale.
\end{enumerate}
